\documentclass[10pt, a4paper]{article}
\usepackage{setspace} 
\usepackage{amsmath, amssymb} 
\usepackage{geometry}
\usepackage{graphicx} 
\usepackage{slashed}
\usepackage{simplewick}
\setlength{\parindent}{0pt}
\usepackage{booktabs} 
\usepackage{float}    



\title{Theory of Solid Note}
\author{Yun-Yen Wu} 
\date{} 

\begin{document}

\maketitle

\tableofcontents

\newpage 

\setcounter{page}{1} 

\section{Introduction}

\section{The Crystal Lattice}

A perfect crystal is defined as an array of atoms arranged in a \textbf{periodic} structure spanning $\mathbb{R}^d$. We assume an idealized infinite model, neglecting surface effects and impurities.

\subsection{Fundamental Definitions}
The physical crystal structure $\mathcal{C}$ is the formal mapping of a mathematical abstraction (the lattice $\mathcal{L}$) and a physical entity (the basis $\mathcal{B}$). 

\paragraph{Mathematical Representation}
If the basis consists of $N$ atoms at positions $\{\mathbf{s}_1, \mathbf{s}_2, \dots, \mathbf{s}_N\}$ relative to a lattice point, the crystal is defined as the set:
\begin{equation}
    \mathcal{C} = \{ \mathbf{R} + \mathbf{s}_i \mid \mathbf{R} \in \mathcal{L}, i = 1, \dots, N \}
\end{equation}
where $\mathcal{L}$ is the Bravais lattice.

\subsection{Real-Space Formulation}

\paragraph{Bravais Lattice} 
A $d$-dimensional Bravais lattice is the set of all position vectors $\mathbf{R}$ generated by a set of linearly independent primitive vectors $\{\mathbf{a}_1, \mathbf{a}_2, \dots, \mathbf{a}_d\}$:
\begin{equation}
    \mathbf{R} = \sum^d_{j=1} n_j \mathbf{a}_j, \quad n_j \in \mathbb{Z}
\end{equation}
The lattice exhibits \textbf{discrete translational symmetry}: the environment is invariant under the transformation $\mathbf{r}' = \mathbf{r} + \mathbf{R}$.

\paragraph{Primitive Unit Cell}
A volume $V_c$ which, when translated by all vectors $\mathbf{R}$ in the Bravais lattice, fills all space without overlapping. It contains exactly one lattice point. In $d$-dimensions, the volume is given by the determinant of the basis matrix $A = [\mathbf{a}_1, \dots, \mathbf{a}_d]$:
\begin{equation}
    V_c = |\det(\mathbf{a}_1, \mathbf{a}_2, \dots, \mathbf{a}_d)|
\end{equation}

\paragraph{Wigner-Seitz Cell}
The unique primitive cell that preserves the full symmetry of the lattice. Geometrically, it is defined as the region of Voronoi proximity: the locus of all points in space that are closer to a chosen lattice point than to any other. Formally, for an origin lattice point $\mathbf{0}$:
\begin{equation}
    \mathcal{W}_s = { \mathbf{r} \in \mathbb{R}^3 \mid \text{dist}(\mathbf{r}, \mathbf{0}) \leq \text{dist}(\mathbf{r}, \mathbf{R}) \quad \forall \mathbf{R} \in L, \mathbf{R} \neq \mathbf{0} }
\end{equation}
It is constructed by taking the intersection of the half-spaces defined by the perpendicular bisector planes between the origin and its nearest neighbors.



\subsection{Unit Cells and Lattice Geometry}

\paragraph{Conventional vs. Primitive Cells}
While a \textbf{primitive cell} is the smallest possible repeating unit (containing $Z=1$ lattice point), it often fails to reflect the full point-group symmetry of the system. We define a \textbf{Conventional Unit Cell} as a larger cell ($Z > 1$) chosen to align with the orthogonal axes of the crystal system (e.g., the cubic cell for FCC and BCC structures).



\paragraph{Coordination Number and Density}
The connectivity of a lattice is characterized by the Coordination Number ($Z_{cn}$), defined as the number of nearest-neighbor points to any given lattice point. For standard 3D cubic systems:

\begin{center}
\begin{tabular}{|l|c|c|c|}
\hline
\textbf{Lattice Type} & \textbf{Lattice Points ($Z$)} & \textbf{Coordination ($Z_{cn}$)} & \textbf{Nearest Neighbor distance ($r$)} \\ \hline
Simple Cubic (SC) & 1 & 6 & $a$ \\ \hline
Body-Centered (BCC) & 2 & 8 & $\sqrt{3}a/2$ \\ \hline
Face-Centered (FCC) & 4 & 12 & $a/\sqrt{2}$ \\ \hline
\end{tabular}
\end{center}

\subsection{Crystallographic Directions and Planes}

To describe directions and orientations within the lattice, we utilize the standard notation of indices based on the primitive or conventional vectors.

\paragraph{Crystal Directions}
A direction is defined by the vector $\mathbf{v} = u\mathbf{a}_1 + v\mathbf{a}_2 + w\mathbf{a}_3$. It is denoted as $[uvw]$, using the smallest integer values possible. Negative components are denoted with a bar, e.g., $[\bar{1}00]$.

\paragraph{Miller Indices (Crystal Planes)}
A set of parallel lattice planes is identified by the indices $(hkl)$. These indices are inversely proportional to the intercepts of the plane with the crystal axes.
\begin{enumerate}
    \item Determine the intercepts $x, y, z$ in units of the lattice vectors.
    \item Take the reciprocals: $1/x, 1/y, 1/z$.
    \item Multiply by a common factor to reduce to the smallest integers $(hkl)$.
\end{enumerate}
Mathematically, the plane $(hkl)$ is perpendicular to the reciprocal lattice vector $\mathbf{G} = h\mathbf{b}_1 + k\mathbf{b}_2 + l\mathbf{b}_3$.

\subsection{Reciprocal Space (Momentum Space)}

\paragraph{Reciprocal Lattice Construction}
The set of all wave vectors $\mathbf{G}$ that yield plane waves with the periodicity of the Bravais lattice. The defining condition is $e^{i \mathbf{G} \cdot \mathbf{R}} = 1$. In 3D, the primitive reciprocal vectors $\{\mathbf{b}_1, \mathbf{b}_2, \mathbf{b}_3\}$ are constructed from the direct lattice vectors $\{\mathbf{a}_i\}$ via:
\begin{equation}
    \mathbf{b}_i = 2\pi \frac{\epsilon_{ijk} \mathbf{a}_j \times \mathbf{a}_k}{\mathbf{a}_1 \cdot (\mathbf{a}_2 \times \mathbf{a}_3)}
\end{equation}
which satisfies the orthogonality condition $\mathbf{b}_i \cdot \mathbf{a}_j = 2\pi \delta_{ij}$. The volume of the reciprocal unit cell is $V_{BZ} = (2\pi)^3 / V_c$.

\paragraph{First Brillouin Zone (FBZ)}
The Wigner-Seitz cell of the reciprocal lattice. It contains all $k$-points that are closer to the origin than to any other reciprocal lattice point. The boundaries of the FBZ are defined by the \textbf{Bragg planes}:
\begin{equation}
    \mathbf{k} \cdot \mathbf{G} = \frac{1}{2} |\mathbf{G}|^2
\end{equation}



\subsection{Symmetry Operations and Group Theory}

A crystal's symmetry is described by its \textbf{Space Group} $\mathcal{G}$, the set of isometries $\{ \alpha | \mathbf{t} \}$ that leave the crystal invariant.

\paragraph{Group Algebra}
The composition of two symmetry operations follows the law:
\begin{equation}
    \{ \alpha_2 | \mathbf{t}_2 \} \{ \alpha_1 | \mathbf{t}_1 \} = \{ \alpha_2 \alpha_1 | \alpha_2 \mathbf{t}_1 + \mathbf{t}_2 \}
\end{equation}
The inverse is given by $\{ \alpha | \mathbf{t} \}^{-1} = \{ \alpha^{-1} | -\alpha^{-1} \mathbf{t} \}$.

\paragraph{Crystallographic Restriction Theorem}
In a periodic lattice, only $n$-fold rotational symmetries with $n \in \{1, 2, 3, 4, 6\}$ are allowed.
Consider a rotation matrix $R_\theta$. Its trace in an orthonormal basis is $1 + 2\cos\theta$. In a basis of primitive lattice vectors, the matrix elements must be integers; thus, $\text{Tr}(R_\theta) = M \in \mathbb{Z}$.
\begin{equation}
    \cos\theta = \frac{M-1}{2} \implies \frac{M-1}{2} \in \left\{ -1, -\frac{1}{2}, 0, \frac{1}{2}, 1 \right\}
\end{equation}
This restricts $\theta$ to $0^\circ, 60^\circ, 90^\circ, 120^\circ, 180^\circ$.

\subsection{Point Group Symmetry (Schönflies Notation)}
Point groups $\mathcal{P} \subset O(3)$ fix at least one point. There are \textbf{32 crystallographic point groups} in 3D.
\begin{itemize}
    \item \textbf{Symmorphic Groups:} Can be expressed as a semi-direct product of the point group and the translation group, $\mathcal{G} = \mathcal{T} \rtimes \mathcal{P}$.
    \item \textbf{Non-Symmorphic Groups:} Contain operations like \textbf{screw axes} (rotation + fractional translation) or \textbf{glide planes} (reflection + fractional translation) that cannot be reduced to a point group operation about any origin.
\end{itemize}

\subsection{Stereographic Projection}
The stereographic projection maps symmetry directions from a unit sphere onto a 2D plane. Mathematically, this is an application of the mapping used in complex analysis to represent the \textbf{Riemann Sphere} on the complex plane. For a point $(X, Y, Z)$ on the sphere, the projection $(x, y)$ onto the equatorial plane from the South Pole $(0,0,-1)$ is:

\begin{equation}
    x = \frac{X}{1+Z}, \quad y = \frac{Y}{1+Z}
\end{equation}

This mapping is \textbf{conformal} (angle-preserving) and maps circles on the sphere to circles or lines on the plane. This property is mathematically essential for preserving the angular relationships between crystal faces, which allows for the direct interpretation of Laue diffraction patterns and symmetry poles.

\subsection{The 7 Crystal Systems and 14 Bravais Lattices}
The classification of crystals follows a symmetry hierarchy: 
\textbf{230 Space Groups} $\rightarrow$ \textbf{32 Point Groups} $\rightarrow$ \textbf{14 Bravais Lattices} $\rightarrow$ \textbf{7 Crystal Systems}. 

\begin{table}[H]
\centering
\caption{Classification of the 7 Crystal Systems}
\begin{tabular}{@{}llllc@{}}
\toprule
\textbf{System} & \textbf{Lattice Types} & \textbf{Axial Relations} & \textbf{Angular Relations} & \textbf{Count} \\ \midrule
Cubic & P, I, F & $a = b = c$ & $\alpha = \beta = \gamma = 90^\circ$ & 3 \\
Tetragonal & P, I & $a = b \neq c$ & $\alpha = \beta = \gamma = 90^\circ$ & 2 \\
Orthorhombic & P, I, F, C & $a \neq b \neq c$ & $\alpha = \beta = \gamma = 90^\circ$ & 4 \\
Monoclinic & P, C & $a \neq b \neq c$ & $\alpha = \gamma = 90^\circ \neq \beta$ & 2 \\
Triclinic & P & $a \neq b \neq c$ & $\alpha \neq \beta \neq \gamma \neq 90^\circ$ & 1 \\
Trigonal & R (Rhombohedral) & $a = b = c$ & $\alpha = \beta = \gamma < 120^\circ, \neq 90^\circ$ & 1 \\
Hexagonal & P & $a = b \neq c$ & $\alpha = \beta = 90^\circ, \gamma = 120^\circ$ & 1 \\ \midrule
\textbf{Total} & & & & \textbf{14} \\ \bottomrule
\end{tabular}
\end{table}



\subsection{Common Crystal Structures}
A crystal structure is the physical realization of a lattice $\mathcal{L}$ with a specific basis $\mathcal{B}$. 

\paragraph{Close-Packed Structures}
Most metals crystallize into high-density structures to minimize internal energy.
\begin{itemize}
    \item \textbf{Face-Centered Cubic (FCC):} A cubic lattice with basis $(0,0,0)$. Coordination number $Z=12$. Packing fraction $\approx 0.74$.
    \item \textbf{Hexagonal Close-Packed (HCP):} Simple Hexagonal lattice with a 2-atom basis: $\mathbf{s}_1 = 0, \mathbf{s}_2 = \frac{2}{3}\mathbf{a}_1 + \frac{1}{3}\mathbf{a}_2 + \frac{1}{2}\mathbf{a}_3$. For ideal packing, $c/a = \sqrt{8/3}$.
\end{itemize}

\paragraph{Covalent and Ionic Structures}
Symmetry is often dictated by bonding geometry rather than just packing efficiency.

\begin{table}[H]
\centering
\caption{Mathematical Basis for Common Structures}
\begin{tabular}{@{}llll@{}}
\toprule
\textbf{Structure} & \textbf{Lattice} & \textbf{Basis Vectors} & \textbf{Examples} \\ \midrule
NaCl & FCC & $(0,0,0), (\frac{1}{2},0,0)$ & LiF, MgO, AgBr \\
CsCl & Simple Cubic & $(0,0,0), (\frac{1}{2},\frac{1}{2},\frac{1}{2})$ & CsI, TlBr, CuZn \\
Diamond & FCC & $(0,0,0), (\frac{1}{4},\frac{1}{4},\frac{1}{4})$ & C, Si, Ge, Sn \\
Zincblende & FCC & $(0,0,0)_A, (\frac{1}{4},\frac{1}{4},\frac{1}{4})_B$ & GaAs, InP, ZnS \\
Perovskite & Simple Cubic & $A(0,0,0), B(\frac{1}{2},\frac{1}{2},\frac{1}{2}), O(0,\frac{1}{2},\frac{1}{2})\text{+}sym$ & $CaTiO_3, BaTiO_3$ \\ \bottomrule
\end{tabular}
\end{table}

\paragraph{Note on Diamond and Zincblende}
Both structures consist of two interpenetrating FCC lattices displaced by one-fourth of the body diagonal. In Diamond, both lattices consist of the same atomic species; in Zincblende, they differ, breaking the inversion symmetry and leading to properties like piezoelectricity in GaAs.

\section{Electrons in a Perioic Potential}

\subsection{Bloch's Theorem}
Bloch's theorem describes the eigenstates of a particle in a periodic potential $V(\mathbf{r} + \mathbf{R}) = V(\mathbf{r})$, where $\mathbf{R}$ is a lattice vector. It states that the one-particle states are given by:
\begin{equation}
    \psi_{n\mathbf{k}}(\mathbf{r}) = u_{n\mathbf{k}}(\mathbf{r}) e^{i\mathbf{k}\cdot\mathbf{r}}
\end{equation}
where $u_{n\mathbf{k}}(\mathbf{r})$ is a periodic function with the same periodicity as the lattice ($u(\mathbf{r} + \mathbf{R}) = u(\mathbf{r})$), and $\mathbf{k}$ is the crystal momentum restricted to the first Brillouin zone.

\subsection{Symmetry and Translation Operators}
Let $\mathbf{R}$ be a lattice vector and $T_\mathbf{R}$ be the translation operator such that $T_\mathbf{R}f(\mathbf{r}) = f(\mathbf{r+R})$. Because the Hamiltonian $H$ is translationally invariant under lattice translations:
\begin{equation}
    [T_\mathbf{R}, H] = 0
\end{equation} 
Consequently, $H$ and the set $\{T_\mathbf{R}\}$ share a simultaneous set of eigenfunctions. Since $T_\mathbf{R}$ is unitary, its eigenvalues are of the form $e^{i\mathbf{k}\cdot\mathbf{R}}$, leading directly to the Bloch condition:
\begin{equation}
    \psi(\mathbf{r} + \mathbf{R}) = e^{i\mathbf{k}\cdot\mathbf{R}} \psi(\mathbf{r})
\end{equation}

\subsection{The $k$-dependent Schrödinger Equation}
The crystal momentum $\hbar \mathbf{k}$ is a quasi-momentum. While not a "true" momentum ($\mathbf{p} \neq \hbar\mathbf{k}$), it is conserved in processes involving the periodic lattice.
By substituting the Bloch form into the Schrödinger equation, we find the equation for the periodic part $u_{n\mathbf{k}}(\mathbf{r})$:
\begin{equation}
    H(\mathbf{k}) u_{n\mathbf{k}}(\mathbf{r}) = \left[ \frac{(\mathbf{p} + \hbar\mathbf{k})^2}{2m} + V(\mathbf{r}) \right] u_{n\mathbf{k}}(\mathbf{r}) = E_{n\mathbf{k}} u_{n\mathbf{k}}(\mathbf{r})
\end{equation}
Expanding the kinetic operator, the effective Hamiltonian is:
\begin{equation}
    H(\mathbf{k}) = \frac{\mathbf{p}^2}{2m} + V(\mathbf{r}) + \frac{\hbar}{m}\mathbf{k}\cdot\mathbf{p} + \frac{\hbar^2k^2}{2m}
\end{equation}
The electron velocity in band $n$ is given by the group velocity:
\begin{equation}
    \mathbf{v}_n(\mathbf{k}) = \frac{1}{\hbar} \nabla_{\mathbf{k}} E_{n\mathbf{k}} = \langle \psi_{n\mathbf{k}} | \frac{\mathbf{p}}{m} | \psi_{n\mathbf{k}} \rangle
\end{equation}

\subsection{Reciprocal Space and the Central Equation}
To solve for the energy bands, we expand the periodic potential and the Bloch function in a Fourier series over reciprocal lattice vectors $\mathbf{G}$:
\begin{align}
    V(\mathbf{r}) &= \sum_{\mathbf{G}} V_{\mathbf{G}} e^{i\mathbf{G} \cdot \mathbf{r}} \\
    \psi_{\mathbf{k}}(\mathbf{r}) &= \sum_{\mathbf{G}} C_{\mathbf{k}-\mathbf{G}} e^{i(\mathbf{k}-\mathbf{G}) \cdot \mathbf{r}}
\end{align}
where $\mathbf{G} \cdot \mathbf{R} = 2\pi n$. Substituting these into the Schrödinger equation yields the \textbf{Central Equation}:
\begin{equation}
    \left[ \frac{\hbar^2}{2m}(\mathbf{k}-\mathbf{G})^2 - E \right] C_{\mathbf{k}-\mathbf{G}} + \sum_{\mathbf{G}'} V_{\mathbf{G}'-\mathbf{G}} C_{\mathbf{k}-\mathbf{G}'} = 0
\end{equation}

\subsection{The Nearly Free Electron Model: Origin of the Band Gap}

To understand the origin of energy gaps, we treat the periodic potential $V(\mathbf{r})$ as a small perturbation to the free electron gas. The Hamiltonian is:
\begin{equation}
    \hat{H} = -\frac{\hbar^2}{2m}\nabla^2 + V(\mathbf{r})
\end{equation}
Since $V(\mathbf{r})$ is periodic, it can be expanded as a Fourier series over the reciprocal lattice vectors $\mathbf{G}$:
\begin{equation}
    V(\mathbf{r}) = \sum_{\mathbf{G}} U_{\mathbf{G}} e^{i\mathbf{G}\cdot\mathbf{r}}
\end{equation}

\subsection{The Central Equation and Gap Opening}
Using Bloch's theorem, the wavefunction $\psi_{\mathbf{k}}(\mathbf{r})$ is expanded in plane waves. At the Brillouin zone boundary where $|\mathbf{k}| \approx |\mathbf{k}-\mathbf{G}|$, the mixing of two states $C_{\mathbf{k}}$ and $C_{\mathbf{k}-\mathbf{G}}$ dominates. The secular equation becomes:
\begin{equation}
    \begin{vmatrix} 
    \lambda_{\mathbf{k}} - \epsilon & U_{\mathbf{G}} \\
    U_{-\mathbf{G}} & \lambda_{\mathbf{k}-\mathbf{G}} - \epsilon 
    \end{vmatrix} = 0
\end{equation}
where $\lambda_{\mathbf{k}} = \frac{\hbar^2 k^2}{2m}$. Solving for $\epsilon$ at the zone boundary ($\lambda_{\mathbf{k}} = \lambda_{\mathbf{k}-\mathbf{G}}$) yields:
\begin{equation}
    \epsilon = \lambda \pm |U_{\mathbf{G}}|
\end{equation}
This demonstrates that the potential $U_{\mathbf{G}}$ lifts the degeneracy at the zone boundary, creating a forbidden energy gap of magnitude $E_g = 2|U_{\mathbf{G}}|$.

\subsection{Electronic Band Structure and Filling}
The classification of solids depends on the occupancy of these bands relative to the \textbf{Fermi level} ($E_F$) and the presence of band overlaps.

\paragraph{Monovalent Metal}
A material with one valence electron per atom and one atom per unit cell. Since each Brillouin zone contains $N$ discrete $\mathbf{k}$-points and can hold $2N$ electrons (due to spin degeneracy), the first conduction band is exactly \textbf{half-filled}.

\paragraph{Even-Valency Metal}
In these materials, the number of electrons is sufficient to fill a band ($2N$). However, if the \textbf{band overlap} is large (i.e., the maximum of a lower band exceeds the minimum of a higher band), electrons spill into the higher band. This creates a partially filled conduction band and a partially empty valence band.

\paragraph{Insulator}
A material with an even number of electrons per unit cell and a large band gap $E_g$ (typically $> 3 \text{ eV}$). At $T = 0 \text{ K}$, the valence band is completely filled and the conduction band is completely empty. The energy $k_B T$ is insufficient to excite carriers across the gap.

\paragraph{Semiconductor}
Technically similar to an insulator but with a small $E_g$ (typically $< 3 \text{ eV}$). At finite temperatures, thermal excitation allows a small number of electrons to reach the conduction band, leaving behind \textbf{holes} in the valence band. Both act as charge carriers.

\paragraph{Semimetal}
An even-valency material where the band overlap is very small. Unlike a metal, the density of states at the Fermi level is low; unlike a semiconductor, there is no true energy gap. The number of electrons equals the number of holes, but both concentrations are small.


\subsection{Reciprocal Space Representations}
Energy bands $E_n(\mathbf{k})$ are periodic in $\mathbf{k}$-space. We visualize this using three schemes:
\begin{itemize}
    \item \textbf{Extended Zone Scheme:} $E(\mathbf{k})$ is plotted as a quasi-free electron parabola, with discontinuities (gaps) at each Bragg plane.
    \item \textbf{Reduced Zone Scheme:} All branches are translated by reciprocal lattice vectors $\mathbf{G}$ into the \textbf{First Brillouin Zone}. This is the standard representation for electronic properties.
    \item \textbf{Repeated Zone Scheme:} The first BZ is repeated periodically. This is useful for visualizing the connectivity of the Fermi surface.
\end{itemize}

\subsection{Higher-Order Brillouin Zones}
Zones are defined by the number of Bragg planes crossed from the origin:
\begin{itemize}
    \item \textbf{1st Brillouin Zone:} The set of points reachable from the origin without crossing any Bragg planes.
    \item \textbf{$n$-th Brillouin Zone:} The set of points reachable from the origin by crossing exactly $n-1$ Bragg planes.
\end{itemize}


\subsection{Fermi Surface and Density of States}
The Fermi surface is defined as the constant-energy manifold in reciprocal space where $E_n(\mathbf{k}) = E_F$. To calculate the Density of States (DOS) per unit volume, we convert the sum over $\mathbf{k}$ into an integral over the Brillouin zone:
\begin{equation}
    D(E) = \frac{g_s}{(2\pi)^d} \sum_n \int_{BZ} d^d k \, \delta (E - E_n(\mathbf{k}))
\end{equation}
By resolving the integral along the direction normal to the constant-energy surface, we obtain the professional form:
\begin{equation}
    D(E) = \frac{g_s}{(2\pi)^d} \sum_n \int_{E_n(\mathbf{k})=E} \frac{dS_n}{|\nabla_{\mathbf{k}} E_n(\mathbf{k})|}
\end{equation}
The points where $\nabla_{\mathbf{k}} E_n(\mathbf{k}) = 0$ are known as \textbf{van Hove singularities}. In 1D, $D(E)$ exhibits square-root divergences, while in 3D, the DOS is continuous but possesses discontinuous derivatives (kinks).

\subsection{Tight-Binding Model}
The Tight-Binding (LCAO) approach assumes that electrons are primarily localized on atomic sites. We define the Bloch basis functions as a linear combination of atomic orbitals $\phi_i$:
\begin{equation}
    \psi_{i\mathbf{k}}(\mathbf{r}) = \frac{1}{\sqrt{N}} \sum_{\mathbf{R}_n} e^{i\mathbf{k} \cdot \mathbf{R}_n} \phi_i(\mathbf{r} - \mathbf{R}_n)
\end{equation}
For a system with multiple orbitals, the eigenvalues $E(\mathbf{k})$ are found by solving the secular determinant:
\begin{equation}
    \det | H_{ij}(\mathbf{k}) - E(\mathbf{k}) S_{ij}(\mathbf{k}) | = 0
\end{equation}
where $H_{ij}(\mathbf{k})$ is the transfer integral matrix and $S_{ij}(\mathbf{k})$ is the overlap matrix. In the \textbf{nearest-neighbor approximation}, we only consider hopping integrals $t = \langle \phi(\mathbf{r}) | \hat{H} | \phi(\mathbf{r}-\mathbf{R}_{nn}) \rangle$.

\paragraph{Wannier Functions}
Wannier functions $w_n(\mathbf{r}-\mathbf{R})$ provide a localized orthonormal basis set. They are related to Bloch functions via a Fourier transform:
\begin{equation}
    w_n(\mathbf{r} - \mathbf{R}) = \frac{V}{(2\pi)^d} \int_{BZ} d^d k \, e^{-i\mathbf{k} \cdot \mathbf{R}} \psi_{n\mathbf{k}}(\mathbf{r})
\end{equation}
This allows for a rigorous mapping between localized real-space physics and delocalized reciprocal-space bands.

\end{document}